\chapter{Penggunaan Dasar}

\section{Mengapa Memantau Penggunaan AI?}

Berlangganan lebih dari satu layanan AI berarti mengelola beberapa jendela waktu
kuota yang berbeda: Claude membatasi pesan per 5 jam dan mingguan, OpenAI Codex
membatasi per 5 jam dan mingguan, sementara OpenCode mencatat biaya token secara
langsung. Tanpa pemantauan, ada tiga masalah yang sering dialami:

\begin{itemize}
  \item \textbf{Kuota habis di tengah pekerjaan.} Claude atau Codex berhenti
        merespons tanpa peringatan sebelumnya, memutus alur kerja saat sedang
        fokus.
  \item \textbf{Biaya tak terduga.} Token yang digunakan melalui OpenCode
        langsung ditagih; tanpa visibilitas, tagihan bulanan bisa mengejutkan.
  \item \textbf{Konteks berpindah-pindah.} Memeriksa kuota berarti membuka
        dasbor web masing-masing penyedia, memecah konsentrasi dari terminal.
\end{itemize}

PakAI menyelesaikan ketiga masalah ini dengan satu perintah di terminal. Status
semua penyedia tersedia setiap saat — di status bar tmux, di panel Waybar, maupun
di dashboard TUI layar penuh — tanpa harus meninggalkan lingkungan kerja.

\begin{tipbox}
  PakAI cocok untuk siapa saja yang menggunakan dua atau lebih layanan AI
  berbayar dan ingin tahu kapan harus beralih penyedia atau mengistirahatkan sesi
  agar kuota tidak terbuang.
\end{tipbox}

\section{Memeriksa Status Penggunaan}

Perintah utama untuk melihat ringkasan penggunaan semua penyedia:

\begin{lstlisting}[style=terminal]
$ pakai status
\end{lstlisting}

\begin{figure}[h]
  \centering
  \includegraphics[width=0.85\linewidth]{pakai-status}
  \caption{Output \texttt{pakai status}: ringkasan semua penyedia aktif}
\end{figure}

Output menampilkan:
\begin{itemize}
  \item Nama penyedia dan label yang dikonfigurasi.
  \item Persentase atau biaya penggunaan per jendela waktu (5j, mingguan, bulanan).
  \item Waktu reset atau sisa waktu jendela.
  \item Jumlah penyedia aktif dan waktu pembaruan terakhir.
\end{itemize}

\subsection{Output JSON}

Tambahkan flag \cmd{--json} untuk mendapatkan output dalam format JSON mentah:

\begin{lstlisting}[style=terminal]
$ pakai status --json
\end{lstlisting}

\section{Dashboard TUI Interaktif}

Dashboard TUI menampilkan tampilan penuh layar dengan pembaruan langsung:

\begin{lstlisting}[style=terminal]
$ pakai dashboard
\end{lstlisting}

\begin{figure}[ht]
  \centering
  \includegraphics[width=0.90\linewidth]{pakai-dashboard}
  \caption{Dashboard TUI PakAI: empat penyedia aktif dengan progress bar
           dan sisa waktu reset per jendela waktu}
\end{figure}

Dashboard diperbarui secara otomatis setiap kali daemon menerima data baru dari
penyedia. Tampilan dibagi per penyedia, masing-masing menampilkan baris untuk
setiap jendela waktu yang aktif.

\subsection{Membaca tampilan dashboard}

Setiap baris di dalam dashboard mengikuti format berikut:

\begin{center}
\begin{tabular}{lp{10cm}}
\toprule
\textbf{Kolom} & \textbf{Keterangan} \\
\midrule
Ikon / kode & Ikon Nerd Font penyedia; fallback dua huruf bila font tidak tersedia. \\[4pt]
Nama penyedia & Label penyedia, termasuk sub-penyedia OpenCode. \\[4pt]
Jendela waktu & \texttt{5h} (5 jam), \texttt{w} (mingguan), \texttt{m} (bulanan). \\[4pt]
Progress bar & Bar 12 segmen menunjukkan proporsi penggunaan saat ini. \\[4pt]
Persentase & Angka persentase atau biaya dalam dolar (\texttt{\$}). \\[4pt]
Waktu reset & Sisa waktu sebelum kuota jendela direset. \\
\bottomrule
\end{tabular}
\end{center}

\subsection{Kode warna}

PakAI menggunakan tiga warna untuk persentase penggunaan:

\begin{itemize}
  \item \textbf{Hijau} — penggunaan di bawah ambang peringatan (aman).
  \item \textbf{Kuning} — mendekati batas; pertimbangkan untuk beralih penyedia
        atau mengurangi frekuensi permintaan.
  \item \textbf{Merah} — mendekati atau melebihi batas; kuota hampir habis.
\end{itemize}

Ambang kuning dan merah dapat dikonfigurasi melalui \cmd{pakai config set}
(lihat Bab~\ref{ch:konfigurasi-lanjutan}).

\subsection{Kapan dashboard berguna}

Dashboard layar penuh paling berguna dalam skenario berikut:

\begin{itemize}
  \item \textbf{Sebelum sesi panjang}: periksa penyedia mana yang masih punya
        kuota cukup sebelum memulai pekerjaan berat.
  \item \textbf{Saat jendela waktu hampir reset}: pantau kolom \textit{resets in}
        untuk memilih waktu terbaik memulai ulang sesi.
  \item \textbf{Membandingkan biaya}: lihat akumulasi biaya dolar OpenCode
        berdampingan dengan kuota pesan Claude dan Codex.
  \item \textbf{Debugging penyedia}: tekan \texttt{d} untuk melihat JSON mentah
        yang dikirim setiap penyedia, berguna saat ada anomali data.
\end{itemize}

\textbf{Tombol keyboard:}
\begin{center}
\begin{tabular}{ll}
\toprule
\textbf{Tombol} & \textbf{Fungsi} \\
\midrule
\texttt{q} & Keluar dari dashboard \\
\texttt{r} & Muat ulang data secara manual \\
\texttt{d} & Aktifkan/nonaktifkan mode debug (tampilkan JSON mentah) \\
\bottomrule
\end{tabular}
\end{center}

\begin{infobox}
  Jika daemon belum berjalan, jalankan \cmd{pakai daemon start} terlebih dahulu
  sebelum membuka dashboard.
\end{infobox}

\section{Output untuk Tmux}

Perintah \cmd{pakai tmux} menghasilkan string ringkas untuk status bar tmux:

\begin{lstlisting}[style=terminal]
$ pakai tmux
  7% |  8% |  $0.00
\end{lstlisting}

Output menggunakan ikon Nerd Font untuk setiap penyedia. Pastikan terminal
menggunakan font yang mendukung Nerd Font.

\section{Output untuk Waybar}

Perintah \cmd{pakai waybar} menghasilkan JSON untuk modul kustom Waybar:

\begin{lstlisting}[style=terminal]
$ pakai waybar
\end{lstlisting}

\begin{figure}[h]
  \centering
  \includegraphics[width=0.85\linewidth]{pakai-waybar}
  \caption{Output JSON dari \texttt{pakai waybar}}
\end{figure}

Format output JSON:
\begin{lstlisting}[style=toml]
{
  "text": "<teks tampilan>",
  "tooltip": "<detail per penyedia>",
  "class": "ok|warning|critical|error",
  "percentage": 0
}
\end{lstlisting}

\section{Men-debug Penyedia}

Untuk memeriksa detail data satu penyedia tertentu:

\begin{lstlisting}[style=terminal]
$ pakai provider debug claude
$ pakai provider debug opencode
$ pakai provider debug openai
\end{lstlisting}

\begin{figure}[h]
  \centering
  \includegraphics[width=0.75\linewidth]{pakai-provider-debug}
  \caption{Output \texttt{pakai provider debug claude}}
\end{figure}

\section{Completion Shell}

Aktifkan pelengkapan otomatis perintah (tab completion) untuk shell Anda:

\begin{lstlisting}[style=terminal]
# Bash
$ pakai completion bash >> ~/.bashrc

# Zsh
$ pakai completion zsh >> ~/.zshrc

# Fish
$ pakai completion fish > ~/.config/fish/completions/pakai.fish
\end{lstlisting}
